\begin{helper}
Let \(\nu\) be a measure on a set \(\Sigma\), let \(C\subseteq \Sigma\) be an
atom cell, and let \(E\subseteq C\).  Suppose \(w\ge0\) and
\[
  \operatorname{ofReal}(w)=\nu(C).
\]
Then \(E\) has a normalized nonnegative real representative relative to the
weight \(w\): there is a real \(p\ge0\) such that \(p=0\) whenever \(w=0\) and
\[
  \operatorname{ofReal}(w p)=\nu(E).
\]
\end{helper}
\begin{proof}
Split into the two possible weight cases.

If \(w=0\), then the cell mass identity gives
\[
\nu(C)=\operatorname{ofReal}(0)=0.
\]
Because \(E\subseteq C\), monotonicity of measure gives
\(\nu(E)\le \nu(C)=0\).  Measures are nonnegative, so \(\nu(E)=0\).  Taking
\(p=0\) gives \(p\ge0\), satisfies the required zero-weight convention, and
\[
\operatorname{ofReal}(w p)=\operatorname{ofReal}(0)=0=\nu(E).
\]

Now suppose \(w\ne0\).  Since \(w\ge0\), this means \(w>0\).  Again
\(E\subseteq C\), so
\[
\nu(E)\le\nu(C)=\operatorname{ofReal}(w).
\]
The right-hand side is finite, hence \(\nu(E)\ne\infty\).  Define
\[
  p=\frac{\nu(E)_{\mathbb R}}{w},
\]
where \(\nu(E)_{\mathbb R}\) denotes the real value \((\nu(E)).\operatorname{toReal}\).
This number is nonnegative because \((\nu(E)).\operatorname{toReal}\ge0\) and
\(w>0\).  The algebra of division by a positive real gives
\[
  w p = w\frac{(\nu(E)).\operatorname{toReal}}{w}
      =(\nu(E)).\operatorname{toReal}.
\]
Since \(\nu(E)\) is finite, the extended-nonnegative-real identity
\[
  \operatorname{ofReal}\bigl((\nu(E)).\operatorname{toReal}\bigr)=\nu(E)
\]
applies.  Substituting the preceding equality gives
\[
\operatorname{ofReal}(w p)=\nu(E).
\]
The implication \(w=0\Rightarrow p=0\) is vacuous in this positive branch.
Thus the displayed \(p\) satisfies all required clauses in every case.
\end{proof}
